\begin{table}[htbp]
\centering
\small
\caption{Source-side Predictive Representativity diagnostics for the HAM10000 internal evaluation ORP.}
\label{tab:source-pr-diagnostics}
\setlength{\tabcolsep}{4pt}
\renewcommand{\arraystretch}{1.08}
\begin{tabular}{lrrrrrrr}
\toprule
\textbf{Metric} & \textbf{$n$ decisions} & \textbf{$n_S$} & \textbf{Mean source} & \textbf{SD seeds} & \textbf{Mean SE} & \textbf{Max half-width} & \textbf{Precision adequate} \\
\midrule
Recall / sensitivity & 25 & 1002 & 0.707 & 0.041 & 0.032 & 0.068 & 25/25 \\
AUC-PR & 25 & 1002 & 0.747 & 0.035 & 0.031 & 0.072 & 25/25 \\
F1-score & 25 & 1002 & 0.683 & 0.025 & 0.026 & 0.055 & 25/25 \\
Precision & 25 & 1002 & 0.663 & 0.047 & 0.032 & 0.067 & 25/25 \\
\bottomrule
\end{tabular}
\begin{flushleft}
\footnotesize
Notes: Each row summarizes 25 source-side estimates, corresponding to five architectures trained under five random seeds. Source uncertainty is computed on the non-augmented HAM10000 internal test ORP for each locked model. The approximate standard error is obtained from the bootstrap confidence interval half-width divided by 1.96. The precision-adequacy count uses a documentation tolerance of CI half-width $\leq 0.075$; this tolerance is not a clinical performance threshold and can be modified in the script. Because the interval TAC/ETC table repeats source estimates across target conditions, only the BOSQUE overall rows are retained for this source-side diagnostic.
\end{flushleft}
\end{table}